\documentclass[11pt,a4paper]{article}
\usepackage[utf8]{inputenc}
\usepackage[T1]{fontenc}
\usepackage[brazil]{babel}
\usepackage[margin=2.5cm]{geometry}
\usepackage{titlesec}
\usepackage{enumitem}
\usepackage{booktabs}
\usepackage{tabularx}
\usepackage{array}
\usepackage{xcolor}
\usepackage{hyperref}
\usepackage{fancyhdr}

\definecolor{tfblue}{RGB}{20,40,70}
\definecolor{tfgold}{RGB}{200,150,40}

\titleformat{\section}{\normalfont\Large\bfseries\color{tfblue}}{\thesection}{0.6em}{}
\titleformat{\subsection}{\normalfont\large\bfseries\color{tfblue}}{\thesubsection}{0.6em}{}

\pagestyle{fancy}
\fancyhf{}
\fancyhead[L]{\small\textit{TrackFleet --- Escopo do Projeto}}
\fancyhead[R]{\small\thepage}
\renewcommand{\headrulewidth}{0.4pt}

\newcolumntype{L}{>{\raggedright\arraybackslash}X}

\begin{document}

\begin{center}
{\LARGE\bfseries\color{tfblue} DOCUMENTO DE ESCOPO DO PROJETO}\\[2pt]
{\Large\bfseries TrackFleet --- Gestão de Rastreador de Automóveis}\\[4pt]
{\small\itshape Disciplina de Gerenciamento de Projetos $\cdot$ Framework PMBOK\textregistered\ 8\textsuperscript{a} Edição $\cdot$ Domínio: Escopo}\\[2pt]
{\small Integrantes: Enzo Allebrand e Kerlon Ribeiro (dois GPs) $\cdot$ 4\textsuperscript{o} semestre $\cdot$ Data: 20/08}
\end{center}

\vspace{2mm}
\noindent\textbf{\color{tfgold}\large Escopo do Projeto}

\vspace{1mm}
\noindent\textit{O ``como'': todo o trabalho necessário --- e apenas o necessário --- para entregar o Escopo do Produto descrito no documento correspondente. É o componente mais importante da linha de base, pois nele reside o valor esperado do projeto.}

\section{Descrição do Trabalho do Projeto}
\begin{itemize}[itemsep=2pt]
  \item Levantar requisitos com a transportadora parceira e mapear a integração com a API do provedor de rastreamento piloto.
  \item Projetar a arquitetura da solução (backend, frontend, banco de dados, integração externa).
  \item Desenvolver o backend: API própria, integração com o rastreador e regras de geração de alertas de manutenção.
  \item Desenvolver o frontend: mapa, painel de indicadores, cadastros e relatórios.
  \item Implementar autenticação, perfis de acesso e as medidas de conformidade com a LGPD.
  \item Testar a solução (testes automatizados e testes com dados reais/simulados do rastreador).
  \item Homologar a plataforma com a transportadora piloto.
  \item Treinar os usuários piloto (gestor e motoristas) e coletar feedback.
  \item Documentar o projeto e encerrar a primeira fase, com registro de lições aprendidas.
\end{itemize}

\section{Estrutura Analítica do Projeto (EAP)}
A EAP decompõe o trabalho total em pacotes menores, atribuíveis e mensuráveis, tornando o acordo sobre o escopo visível a todos os envolvidos.

\begin{enumerate}[label=\textbf{\arabic*.}, itemsep=3pt]
  \item \textbf{TrackFleet}
  \begin{enumerate}[label=\textbf{\arabic{enumi}.\arabic*}, itemsep=2pt]
    \item Gestão do Projeto
    \begin{enumerate}[label=\arabic{enumi}.\arabic{enumii}.\arabic*, itemsep=1pt]
      \item Termo de abertura e governança
      \item Planejamento e acompanhamento
    \end{enumerate}
    \item Levantamento e Análise de Requisitos
    \begin{enumerate}[label=\arabic{enumi}.\arabic{enumii}.\arabic*, itemsep=1pt]
      \item Entrevistas com a transportadora parceira
      \item Especificação de requisitos funcionais e não funcionais
    \end{enumerate}
    \item Desenvolvimento
    \begin{enumerate}[label=\arabic{enumi}.\arabic{enumii}.\arabic*, itemsep=1pt]
      \item Backend e integração com a API de rastreamento
      \item Frontend (mapa, painel, cadastros)
      \item Autenticação e conformidade com a LGPD
    \end{enumerate}
    \item Testes e Homologação
    \begin{enumerate}[label=\arabic{enumi}.\arabic{enumii}.\arabic*, itemsep=1pt]
      \item Testes automatizados
      \item Piloto com a transportadora parceira
    \end{enumerate}
    \item Treinamento e Encerramento
    \begin{enumerate}[label=\arabic{enumi}.\arabic{enumii}.\arabic*, itemsep=1pt]
      \item Treinamento dos usuários piloto
      \item Documentação final e lições aprendidas
    \end{enumerate}
  \end{enumerate}
\end{enumerate}

\section{Dicionário da EAP (resumido)}
\begin{center}
\begin{tabularx}{\textwidth}{l L L}
\toprule
\textbf{Pacote} & \textbf{Descrição} & \textbf{Critério de aceite} \\
\midrule
1.3.1 Backend + integração & API própria consumindo dados do rastreador piloto e gerando alertas & Dados do rastreador chegam à plataforma e alertas são disparados corretamente \\
1.3.2 Frontend & Mapa, painel e cadastros funcionais & Gestor localiza um veículo e visualiza alertas em até 3 cliques \\
1.4.2 Piloto & Uso real da plataforma pela transportadora parceira por um período de teste & Transportadora usa a plataforma e valida os indicadores gerados \\
\bottomrule
\end{tabularx}
\end{center}

\section{Fora do Escopo do Projeto (Exclusões)}
\begin{itemize}[itemsep=2pt]
  \item Compra, instalação física ou manutenção de rastreadores GPS.
  \item Suporte comercial e operação em produção após o encerramento do piloto acadêmico.
  \item Integração com mais de um provedor de hardware de rastreamento.
  \item Certificações formais de segurança ou auditorias externas de compliance além da adequação básica à LGPD.
\end{itemize}

\section{Linha de Base do Escopo e Controle de Mudanças}
A linha de base do escopo deste documento é aprovada pelo patrocinador (Diretoria da transportadora parceira), conforme definido no Documento de Governança. Qualquer mudança de escopo que altere objetivo, prazo ou entregáveis essenciais segue o mesmo fluxo de escalonamento já definido: sem consenso entre os dois GPs, ou se a mudança afetar objetivo/prazo/valor esperado, a decisão sobe ao patrocinador e é registrada no log de decisões do projeto. Isso evita o \textit{scope creep} --- o crescimento não controlado do escopo por pequenos acréscimos aceitos ``por fora'' do processo.

\section{Indicadores de Escopo}
\begin{center}
\begin{tabularx}{\textwidth}{L l l}
\toprule
\textbf{Indicador} & \textbf{Fórmula} & \textbf{Leitura} \\
\midrule
Precisão da Definição & Itens entregues corretamente $\div$ Itens planejados $\times$ 100 & Maior é melhor \\
Scope Creep & Entregas não planejadas $\div$ Total de entregas $\times$ 100 & Menor é melhor \\
Estabilidade de Requisitos & Requisitos não alterados $\div$ Total de requisitos $\times$ 100 & Maior é melhor \\
\bottomrule
\end{tabularx}
\end{center}

\noindent Esses indicadores devem ser acompanhados a cada check-in do projeto (conforme a cadência definida na Governança), para que desvios de escopo sejam percebidos antes da entrega final, e não durante ela.

\end{document}
